\chapter{没有指挥的机器}
% 完整版：chapters/02_glutcomputer.tex

你的计算机里有一位指挥——时钟发生器。它每秒挥棒几十亿次，电路里所有元件齐步走，同时翻转。大脑里没有这样的指挥。每个神经元各干各的：信号什么时候到就什么时候到，什么时候发就什么时候发。我们来看看，如果只用\nt{GLUT}神经元——清一色的“加号”——这样一台机器能算些什么。

\section*{漏水的桶}

不妨把神经元想象成一只底部有洞的桶。每来一个脉冲，就是倒进一勺水。水从洞里一点点漏掉。如果一勺接一勺来得很勤，水位来得及升到桶沿——桶就翻倒：神经元咔嗒一声，桶空了，一切从头开始。如果水来得稀稀拉拉，水就先漏光了，什么也不会发生。

\pencilfig{2}{\pencilart{2}}{神经元是一只漏水的桶：输入的水滴把它灌满，水又不断漏走；水位到了刻线——咔嗒。}

这个洞正是神经元和逻辑门的主要区别。神经元不会永远记住来过的脉冲，只记得几毫秒。落在这个时间窗内的一切会相加；更早的就忘了。

\section*{“或”与“与”}

如果一勺水就能让桶翻倒，神经元对任何一个输入都会响应。这就是“或”。如果需要两勺，事情就有意思了：没有指挥，“两个都到了吗？”这个问题毫无意义，除非先说清楚——\emph{在多长的时间之内}。只有两个脉冲几乎同时到达，趁第一勺水还没漏掉，神经元才会触发。这就是“时间上的与”，即符合检测器。有些神经元的时间窗是几毫秒，听觉神经元的则不到一毫秒。

\section*{光有“加号”做不到什么}

试着搭一个“非”——输入沉默时它反而工作的神经元。做不到：沉默不会往桶里倒水。这不只对单个神经元成立，对任何只由“加号”组成的网络都成立：把输入加强，网络里没有一个神经元会因此变得更安静。由此得出最麻烦的一条性质：\emph{不能用活动去熄灭活动}。火一旦烧旺，只有一个办法能让它停下——别再添柴。

大自然找到了一条绕行的路。在许多感官里，同一个刺激由两组神经元来传递。眼睛里有一类细胞报告“变亮了”，另一类报告“变暗了”。如果你用的不是一根导线，而是一对“是”和“否”导线，“非”就白白得到了：把两根导线对调即可。用这样成对的导线，只靠“加号”就能搭出任何逻辑电路——甚至还能知道计算已经结束：每一对里恰好亮了一根导线。工程师就是这样搭建异步电路的，它们在任何延迟下都能正确工作。

\section*{用延迟计算时间}

没有指挥，也就没有公共时钟，但时间照样可以算出来。侧面传来的声音到达近处耳朵的时间比到达远处耳朵要早，早出几分之一毫秒。从左耳引一根导线，沿着一排符合检测器从左往右走；从右耳引一根，从右往左走。两个信号相向而行，在两者同时到达的那个点相遇。正是那里的检测器被触发。时间差变成了被触发神经元的\emph{编号}，也就是声音的方向。精度可达十几微秒，尽管单个神经元远没有这么准：靠的是检测器数量多。

\pencilfig{102}{%
% delay line: the left-ear wire runs right along the top, the right-ear wire runs left along the bottom
\draw[pen=pcBlue] (0,0.6) -- (5.6,0.6); \draw[pen=pcBlue] (6.0,-0.6) -- (0.4,-0.6);
\foreach \i in {1,...,7}{
  \pgfmathsetmacro{\x}{0.75*\i}
  \draw[pen=pcBlue] (\x,0.6) -- (\x,0.24); \draw[pen=pcBlue] (\x,-0.6) -- (\x,-0.24);
  \ifnum\i=5 \draw[dotpen=pcAmber, wash=pcAmber, line width=1.1pt] (\x,0) circle (0.2);
  \else \draw[dotpen=pcBlue, wash=pcBlue] (\x,0) circle (0.2); \fi}
\draw[arr=pcBlue] (0.95,0.78) -- (1.75,0.78); \draw[arr=pcBlue] (5.3,-0.78) -- (4.5,-0.78);
\node[lbl, anchor=east] at (-0.05,0.6) {左耳};
\node[lbl, anchor=west] at (6.05,-0.6) {右耳};
\node[lbl, text=pcAmber!70!black, anchor=south] at (3.75,0.95) {在这里相遇};
\draw[arr=pcAmber!70!black] (3.75,0.95) -- (3.75,0.68);
% sound source on the left
\draw[penlite] (-1.25,-0.25) -- (-1.05,-0.25) -- (-0.8,-0.45) -- (-0.8,0.05) -- (-1.05,-0.15) -- (-1.25,-0.15) -- cycle;
\foreach \r in {0.2,0.35}{\draw[penlite] ($(-0.75,-0.2)+(-40:\r)$) arc (-40:40:\r);}
\node[lbl, anchor=north] at (-0.95,-0.5) {声音在左};
}{左边来的声音先到左耳，它的信号跑得更远，于是相遇点落在中点右侧。被触发神经元的编号就是声音的方向。}

\section*{记忆是一堆篝火}

取一群彼此强烈兴奋的神经元。用一个短信号点燃它，它就能自我维持，像一堆自己给自己添柴的篝火。信号结束了，这群神经元却还在燃烧——它记得信号来过。这就是一个存储单元。如果信号太短，篝火来不及烧旺，就熄灭了。

可麻烦在于：这样的记忆怎么擦掉？靠活动——办不到。只能等神经元累了，篝火自己烧尽。

\section*{由事件启动的定时器}

把这样的神经元群连成一条链：第一群点燃第二群，第二群点燃第三群。推一下第一群——一道波就沿链条跑下去，像推倒一排多米诺骨牌。倒下的是第几块骨牌，就说明从推的那一刻起过去了多久。这不是一直嘀嗒走个不停的钟，而是一个由事件启动、只运行一次的定时器。鸣禽唱歌时，某个脑区的神经元严格依次触发——很像这样的链条。

总之，光用“加号”就能做很多事。还缺三样：从许多选项中选出一个，按命令擦除记忆，以及防止网络烧成一片大火。这三样都要靠“减号”——那是下一章的内容。

\fullversion{2}
